
\documentclass{article}
\usepackage{colortbl}
\usepackage{makecell}
\usepackage{multirow}
\usepackage{supertabular}

\begin{document}

\newcounter{utterance}

\centering \large Interaction Transcript for game `dond', experiment `semi', episode 2 with deepseek{-}v4{-}pro.
\vspace{24pt}

{ \footnotesize  \setcounter{utterance}{1}
\setlength{\tabcolsep}{0pt}
\begin{supertabular}{c@{$\;$}|p{.15\linewidth}@{}p{.15\linewidth}p{.15\linewidth}p{.15\linewidth}p{.15\linewidth}p{.15\linewidth}}
    \# & $\;$A & \multicolumn{4}{c}{Game Master} & $\;\:$B\\
    \hline

    \theutterance \stepcounter{utterance}  
    & & \multicolumn{4}{p{0.6\linewidth}}{
        \cellcolor[rgb]{0.9,0.9,0.9}{
            \makecell[{{p{\linewidth}}}]{
                \texttt{\tiny{[P1$\langle$GM]}}
                \texttt{Igrate kooperativnu igru pregovaranja u kojoj se trebate dogovoriti s drugim igračem o tome kako podijeliti skup stavki između vas.} \\
\\ 
\texttt{Pravila:} \\
\texttt{(a) Vi i drugi igrač dobivate skup stavki. Svakom od vas također se daje tajna funkcija vrijednosti koja predstavlja koliko cijenite svaku vrstu predmeta.} \\
\texttt{(b) Razmjenjujete poruke s drugim igračem kako biste se dogovorili o tome tko dobiva koje stavke. Svaki može poslati maksimalno 5 poruka ili prekinuti igru ranije dajući tajni prijedlog u bilo kojem trenutku.} \\
\texttt{(c) Od vas se traži da podnese tajni prijedlog koji označava stavke koje želite, formatirane u uglatim zagradama kako slijedi: "{[}Prijedlog: <broj> <naziv stavke>, <broj> <naziv stavke>, <...>{]}"} \\
\texttt{(d) Ako su vaši prijedlozi komplementarni, tj. ima dovoljno stavki za ispunjavanje oba prijedloga, svaki igrač dobiva rezultat na temelju zbroja vrijednosti stavki koje je dobio. U suprotnom, obojica dobivate nula bodova.} \\
\texttt{(e) Vaš cilj je maksimizirati zbroj vašeg rezultata i rezultata drugog igrača.} \\
\\ 
\texttt{Počnimo.} \\
\\ 
\texttt{Skup dostupnih stavki je:} \\
\\ 
\texttt{1 torba, 3 lampe, 1 cvijet} \\
\\ 
\texttt{Vaša tajna funkcija vrijednosti je:} \\
\\ 
\texttt{torba: 3, lampa: 0, cvijet: 7} \\
\\ 
\texttt{VAŽNO: Vaše poruke, osim ako nisu tajni prijedlog, šalju se izravno drugom igraču, pa ne uključujte nikakav odgovor na pravila niti tekst koji najavljuje vašu poruku. Za davanje tajnog prijedloga koristite navedeni format. Ne koristite uglate zagrade pri komuniciranju s drugim igračem ili će se tumačiti kao vaš tajni prijedlog.} \\
            }
        }
    }
    & \\ \\

    \theutterance \stepcounter{utterance}  
    & \multicolumn{4}{p{0.6\linewidth}}{
        \cellcolor[rgb]{0.9,0.9,0.9}{
            \makecell[{{p{\linewidth}}}]{
                \texttt{\tiny{[P1$\rangle$GM]}}
                \texttt{Shvaćam situaciju. Imam 5 poruka za komunikaciju prije davanja prijedloga. Krenimo s pregovorima.} \\
\\ 
\texttt{Pozdrav! Vidim da imamo torbu, tri lampe i jedan cvijet za podijeliti. Ja osobno najviše cijenim cvijet, dok su mi lampe nebitne. Što ti najviše vrijedi od ovih predmeta?} \\
            }
        }
    }
    & & \\ \\

    \theutterance \stepcounter{utterance}  
    & & & \multicolumn{2}{p{0.3\linewidth}}{
        \cellcolor[rgb]{0.9,0.9,0.9}{
            \makecell[{{p{\linewidth}}}]{
                \texttt{\tiny{[GM$|$GM]}}
                \texttt{continue} \\
            }
        }
    }
    & & \\ \\

    \theutterance \stepcounter{utterance}  
    & & & \multicolumn{2}{p{0.3\linewidth}}{
        \cellcolor[rgb]{0.9,0.9,0.9}{
            \makecell[{{p{\linewidth}}}]{
                \texttt{\tiny{[GM$|$GM]}}
                \texttt{Shvaćam situaciju. Imam 5 poruka za komunikaciju prije davanja prijedloga. Krenimo s pregovorima.} \\
\\ 
\texttt{Pozdrav! Vidim da imamo torbu, tri lampe i jedan cvijet za podijeliti. Ja osobno najviše cijenim cvijet, dok su mi lampe nebitne. Što ti najviše vrijedi od ovih predmeta?} \\
            }
        }
    }
    & & \\ \\

    \theutterance \stepcounter{utterance}  
    & & \multicolumn{4}{p{0.6\linewidth}}{
        \cellcolor[rgb]{0.9,0.9,0.9}{
            \makecell[{{p{\linewidth}}}]{
                \texttt{\tiny{[GM$\rangle$P2]}}
                \texttt{Igrate kooperativnu igru pregovaranja u kojoj se trebate dogovoriti s drugim igračem o tome kako podijeliti skup stavki između vas.} \\
\\ 
\texttt{Pravila:} \\
\texttt{(a) Vi i drugi igrač dobivate skup stavki. Svakom od vas također se daje tajna funkcija vrijednosti koja predstavlja koliko cijenite svaku vrstu predmeta.} \\
\texttt{(b) Razmjenjujete poruke s drugim igračem kako biste se dogovorili o tome tko dobiva koje stavke. Svaki može poslati maksimalno 5 poruka ili prekinuti igru ranije dajući tajni prijedlog u bilo kojem trenutku.} \\
\texttt{(c) Od vas se traži da podnese tajni prijedlog koji označava stavke koje želite, formatirane u uglatim zagradama kako slijedi: "{[}Prijedlog: <broj> <naziv stavke>, <broj> <naziv stavke>, <...>{]}"} \\
\texttt{(d) Ako su vaši prijedlozi komplementarni, tj. ima dovoljno stavki za ispunjavanje oba prijedloga, svaki igrač dobiva rezultat na temelju zbroja vrijednosti stavki koje je dobio. U suprotnom, obojica dobivate nula bodova.} \\
\texttt{(e) Vaš cilj je maksimizirati zbroj vašeg rezultata i rezultata drugog igrača.} \\
\\ 
\texttt{Počnimo.} \\
\\ 
\texttt{Skup dostupnih stavki je:} \\
\\ 
\texttt{1 torba, 3 lampe, 1 cvijet} \\
\\ 
\texttt{Vaša tajna funkcija vrijednosti je:} \\
\\ 
\texttt{torba: 1, lampa: 2, cvijet: 3} \\
\\ 
\texttt{VAŽNO: Vaše poruke, osim ako nisu tajni prijedlog, šalju se izravno drugom igraču, pa ne uključujte nikakav odgovor na pravila niti tekst koji najavljuje vašu poruku. Za davanje tajnog prijedloga koristite navedeni format. Ne koristite uglate zagrade pri komuniciranju s drugim igračem ili će se tumačiti kao vaš tajni prijedlog.} \\
\\ 
\texttt{Shvaćam situaciju. Imam 5 poruka za komunikaciju prije davanja prijedloga. Krenimo s pregovorima.} \\
\\ 
\texttt{Pozdrav! Vidim da imamo torbu, tri lampe i jedan cvijet za podijeliti. Ja osobno najviše cijenim cvijet, dok su mi lampe nebitne. Što ti najviše vrijedi od ovih predmeta?} \\
            }
        }
    }
    & \\ \\

    \theutterance \stepcounter{utterance}  
    & & & \multicolumn{4}{p{0.6\linewidth}}{
        \cellcolor[rgb]{0.9,0.9,0.9}{
            \makecell[{{p{\linewidth}}}]{
                \texttt{\tiny{[GM$\langle$P2]}}
                \texttt{Pozdrav! Ja najviše cijenim lampe {-} one su mi najvrjednije. Torba i cvijet su mi manje važni. Možemo li se dogovoriti da ja uzmem lampe, a ti dobiješ torbu i cvijet? Tako bismo oboje dobili ono što nam je važno.} \\
            }
        }
    }
     \\ \\

    \theutterance \stepcounter{utterance}  
    & & & \multicolumn{2}{p{0.3\linewidth}}{
        \cellcolor[rgb]{0.9,0.9,0.9}{
            \makecell[{{p{\linewidth}}}]{
                \texttt{\tiny{[GM$|$GM]}}
                \texttt{continue} \\
            }
        }
    }
    & & \\ \\

    \theutterance \stepcounter{utterance}  
    & & & \multicolumn{2}{p{0.3\linewidth}}{
        \cellcolor[rgb]{0.9,0.9,0.9}{
            \makecell[{{p{\linewidth}}}]{
                \texttt{\tiny{[GM$|$GM]}}
                \texttt{Pozdrav! Ja najviše cijenim lampe {-} one su mi najvrjednije. Torba i cvijet su mi manje važni. Možemo li se dogovoriti da ja uzmem lampe, a ti dobiješ torbu i cvijet? Tako bismo oboje dobili ono što nam je važno.} \\
            }
        }
    }
    & & \\ \\

    \theutterance \stepcounter{utterance}  
    & & \multicolumn{4}{p{0.6\linewidth}}{
        \cellcolor[rgb]{0.9,0.9,0.9}{
            \makecell[{{p{\linewidth}}}]{
                \texttt{\tiny{[P1$\langle$GM]}}
                \texttt{Pozdrav! Ja najviše cijenim lampe {-} one su mi najvrjednije. Torba i cvijet su mi manje važni. Možemo li se dogovoriti da ja uzmem lampe, a ti dobiješ torbu i cvijet? Tako bismo oboje dobili ono što nam je važno.} \\
            }
        }
    }
    & \\ \\

    \theutterance \stepcounter{utterance}  
    & \multicolumn{4}{p{0.6\linewidth}}{
        \cellcolor[rgb]{0.9,0.9,0.9}{
            \makecell[{{p{\linewidth}}}]{
                \texttt{\tiny{[P1$\rangle$GM]}}
                \texttt{To zvuči kao dobar prijedlog! Ako ti uzmeš sve tri lampe, a ja uzmem torbu i cvijet, oboje dobivamo ono što nam je najvrjednije. Samo da potvrdim {-} jesu li ti lampe toliko važne da bi ti ta podjela donijela maksimalan rezultat? Meni bi torba i cvijet bili odlični.} \\
            }
        }
    }
    & & \\ \\

    \theutterance \stepcounter{utterance}  
    & & & \multicolumn{2}{p{0.3\linewidth}}{
        \cellcolor[rgb]{0.9,0.9,0.9}{
            \makecell[{{p{\linewidth}}}]{
                \texttt{\tiny{[GM$|$GM]}}
                \texttt{continue} \\
            }
        }
    }
    & & \\ \\

    \theutterance \stepcounter{utterance}  
    & & & \multicolumn{2}{p{0.3\linewidth}}{
        \cellcolor[rgb]{0.9,0.9,0.9}{
            \makecell[{{p{\linewidth}}}]{
                \texttt{\tiny{[GM$|$GM]}}
                \texttt{To zvuči kao dobar prijedlog! Ako ti uzmeš sve tri lampe, a ja uzmem torbu i cvijet, oboje dobivamo ono što nam je najvrjednije. Samo da potvrdim {-} jesu li ti lampe toliko važne da bi ti ta podjela donijela maksimalan rezultat? Meni bi torba i cvijet bili odlični.} \\
            }
        }
    }
    & & \\ \\

    \theutterance \stepcounter{utterance}  
    & & \multicolumn{4}{p{0.6\linewidth}}{
        \cellcolor[rgb]{0.9,0.9,0.9}{
            \makecell[{{p{\linewidth}}}]{
                \texttt{\tiny{[GM$\rangle$P2]}}
                \texttt{To zvuči kao dobar prijedlog! Ako ti uzmeš sve tri lampe, a ja uzmem torbu i cvijet, oboje dobivamo ono što nam je najvrjednije. Samo da potvrdim {-} jesu li ti lampe toliko važne da bi ti ta podjela donijela maksimalan rezultat? Meni bi torba i cvijet bili odlični.} \\
            }
        }
    }
    & \\ \\

    \theutterance \stepcounter{utterance}  
    & & & \multicolumn{4}{p{0.6\linewidth}}{
        \cellcolor[rgb]{0.9,0.9,0.9}{
            \makecell[{{p{\linewidth}}}]{
                \texttt{\tiny{[GM$\langle$P2]}}
                \texttt{Da, potpuno! Lampe su mi daleko najvrjednije, tako da bih s tri lampe bio maksimalno zadovoljan. Torba i cvijet za tebe, lampe za mene {-} to je win{-}win podjela koja maksimizira naše zajedničke bodove. Spremna sam dati prijedlog. Hoćemo li oboje sada poslati tajne prijedloge?} \\
            }
        }
    }
     \\ \\

    \theutterance \stepcounter{utterance}  
    & & & \multicolumn{2}{p{0.3\linewidth}}{
        \cellcolor[rgb]{0.9,0.9,0.9}{
            \makecell[{{p{\linewidth}}}]{
                \texttt{\tiny{[GM$|$GM]}}
                \texttt{continue} \\
            }
        }
    }
    & & \\ \\

    \theutterance \stepcounter{utterance}  
    & & & \multicolumn{2}{p{0.3\linewidth}}{
        \cellcolor[rgb]{0.9,0.9,0.9}{
            \makecell[{{p{\linewidth}}}]{
                \texttt{\tiny{[GM$|$GM]}}
                \texttt{Da, potpuno! Lampe su mi daleko najvrjednije, tako da bih s tri lampe bio maksimalno zadovoljan. Torba i cvijet za tebe, lampe za mene {-} to je win{-}win podjela koja maksimizira naše zajedničke bodove. Spremna sam dati prijedlog. Hoćemo li oboje sada poslati tajne prijedloge?} \\
            }
        }
    }
    & & \\ \\

    \theutterance \stepcounter{utterance}  
    & & \multicolumn{4}{p{0.6\linewidth}}{
        \cellcolor[rgb]{0.9,0.9,0.9}{
            \makecell[{{p{\linewidth}}}]{
                \texttt{\tiny{[P1$\langle$GM]}}
                \texttt{Da, potpuno! Lampe su mi daleko najvrjednije, tako da bih s tri lampe bio maksimalno zadovoljan. Torba i cvijet za tebe, lampe za mene {-} to je win{-}win podjela koja maksimizira naše zajedničke bodove. Spremna sam dati prijedlog. Hoćemo li oboje sada poslati tajne prijedloge?} \\
            }
        }
    }
    & \\ \\

    \theutterance \stepcounter{utterance}  
    & \multicolumn{4}{p{0.6\linewidth}}{
        \cellcolor[rgb]{0.9,0.9,0.9}{
            \makecell[{{p{\linewidth}}}]{
                \texttt{\tiny{[P1$\rangle$GM]}}
                \texttt{Odlično, dogovoreno! Šaljemo prijedloge {-} ja tražim torbu i cvijet, ti tražiš tri lampe. Dajmo oboje svoje prijedloge sada.} \\
\\ 
\texttt{{[}Prijedlog: 1 torba, 1 cvijet{]}} \\
            }
        }
    }
    & & \\ \\

    \theutterance \stepcounter{utterance}  
    & & & \multicolumn{2}{p{0.3\linewidth}}{
        \cellcolor[rgb]{0.9,0.9,0.9}{
            \makecell[{{p{\linewidth}}}]{
                \texttt{\tiny{[GM$|$GM]}}
                \texttt{proposal} \\
            }
        }
    }
    & & \\ \\

    \theutterance \stepcounter{utterance}  
    & & & \multicolumn{2}{p{0.3\linewidth}}{
        \cellcolor[rgb]{0.9,0.9,0.9}{
            \makecell[{{p{\linewidth}}}]{
                \texttt{\tiny{[GM$|$GM]}}
                [1, 0, 1]
            }
        }
    }
    & & \\ \\

    \theutterance \stepcounter{utterance}  
    & & \multicolumn{4}{p{0.6\linewidth}}{
        \cellcolor[rgb]{0.9,0.9,0.9}{
            \makecell[{{p{\linewidth}}}]{
                \texttt{\tiny{[GM$\rangle$P2]}}
                \texttt{Drugi igrač je poslao tajni prijedlog. Sada je na tebi red da pošalješ vlastiti prijedlog.} \\
            }
        }
    }
    & \\ \\

    \theutterance \stepcounter{utterance}  
    & & & \multicolumn{4}{p{0.6\linewidth}}{
        \cellcolor[rgb]{0.9,0.9,0.9}{
            \makecell[{{p{\linewidth}}}]{
                \texttt{\tiny{[GM$\langle$P2]}}
                \texttt{{[}Prijedlog: 1 torba, 1 cvijet{]}} \\
            }
        }
    }
     \\ \\

    \theutterance \stepcounter{utterance}  
    & & & \multicolumn{2}{p{0.3\linewidth}}{
        \cellcolor[rgb]{0.9,0.9,0.9}{
            \makecell[{{p{\linewidth}}}]{
                \texttt{\tiny{[GM$|$GM]}}
                \texttt{proposal} \\
            }
        }
    }
    & & \\ \\

    \theutterance \stepcounter{utterance}  
    & & & \multicolumn{2}{p{0.3\linewidth}}{
        \cellcolor[rgb]{0.9,0.9,0.9}{
            \makecell[{{p{\linewidth}}}]{
                \texttt{\tiny{[GM$|$GM]}}
                [1, 0, 1]
            }
        }
    }
    & & \\ \\

    \theutterance \stepcounter{utterance}  
    & & & \multicolumn{2}{p{0.3\linewidth}}{
        \cellcolor[rgb]{0.9,0.9,0.9}{
            \makecell[{{p{\linewidth}}}]{
                \texttt{\tiny{[GM$|$GM]}}
                [[1, 0, 1], [1, 0, 1]]
            }
        }
    }
    & & \\ \\

\end{supertabular}
}

\end{document}
